% !TeX root = ../main.tex
\documentclass[../main.tex]{subfiles}

\begin{document}

\section{Implementación}

La sección anterior ha descrito el diseño del sistema desde una perspectiva arquitectónica: bloques funcionales, contratos entre módulos y decisiones de diseño adoptadas para transformar una señal EEG real en una salida MIDI física. En esta sección se detalla cómo se ha materializado esa arquitectura en código, indicando qué responsabilidades asume cada parte del sistema y cómo circulan los datos entre ellas.

El criterio seguido en la implementación es mantener separadas las capas de adquisición, transporte, procesamiento, sonificación y salida MIDI. De esta forma, el firmware del microcontrolador se concentra en las tareas de tiempo real asociadas al ADS1299, mientras que el entorno Linux/Python de la Arduino UNO Q recibe los bloques EEG ya escalados, calcula características temporales y espectrales, evalúa la calidad de la señal y genera los eventos musicales. Esta separación evita que la lógica musical dependa directamente de detalles de bajo nivel del convertidor analógico-digital y permite analizar cada etapa de forma independiente.

\subsection{Implementación del firmware de adquisición}

El firmware de adquisición constituye la primera capa ejecutable del sistema. Su papel es transformar la salida digital del ADS1299 en una señal EEG preparada para el backend Python. Para ello configura el convertidor, espera a que haya una nueva conversión disponible, lee el frame por SPI, reconstruye las muestras, las escala a unidades físicas, aplica un filtrado digital básico y agrupa los datos en bloques mediante el evento \texttt{eeg\_block\_uV}.

Dentro del flujo general del sistema, esta capa se encarga de las tareas más cercanas al hardware y al tiempo real. El microcontrolador asegura que las muestras lleguen a Python con una cadencia coherente, en microvoltios y con una estructura de bloque estable. A partir de esos bloques, el backend Python realiza el almacenamiento temporal, el análisis espectral, la evaluación de calidad y la generación de parámetros para la sonificación.

La Figura~\ref{fig:impl_firmware_flujo_adquisicion} resume esta cadena de transformación, desde el frame nativo del ADS1299 hasta el bloque \texttt{eeg\_block\_uV} consumido por el backend Python.

\begin{figure}[H]
    \centering
    \IfFileExists{images/uml/impl_firmware_flujo_adquisicion.pdf}{%
        \includegraphics[width=0.92\textwidth]{images/uml/impl_firmware_flujo_adquisicion.pdf}
    }{%
        \fbox{\parbox{0.9\textwidth}{\centering Diagrama pendiente de exportar desde Mermaid: \texttt{images/uml/impl\_firmware\_flujo\_adquisicion.mmd}}}%
    }
    \caption{Flujo de adquisición implementado en el firmware. El microcontrolador lee frames RDATAC del ADS1299, valida el campo de estado, reconstruye las muestras de 24 bits, las escala, aplica filtrado digital y publica bloques de 8 muestras hacia Linux.}
    \label{fig:impl_firmware_flujo_adquisicion}
\end{figure}

La implementación se reparte en varios archivos. \texttt{sketch.ino} actúa como programa principal: inicializa Bridge, Monitor, ADS1299, filtros y handlers de comunicación, y ejecuta el bucle de adquisición. \texttt{ADS1299\_SafeSPI} encapsula el transporte SPI de bajo nivel; \texttt{ADS1299Plus} implementa la secuencia de arranque, los comandos, el acceso a registros y la lectura de frames; \texttt{ADS1299\_Registers.h} recoge comandos y direcciones con nombres legibles; y \texttt{filters.h} y \texttt{streaming.h} contienen, respectivamente, el acondicionamiento digital y la publicación de bloques hacia Linux.

\subsubsection{Configuración del ADS1299}

El ADS1299-4 es un convertidor analógico-digital de cuatro canales con 24 bits de resolución, orientado a señales EEG y biopotenciales. Integra amplificador de ganancia programable, referencia de tensión, oscilador y filtro digital de decimación. En este proyecto trabaja a \SI{250}{\hertz}, con ganancia 24 y referencia interna de \SI{4.5}{\volt}.

La configuración del ADS1299 se realiza en dos etapas. Primero, el firmware deja el convertidor en un estado conocido mediante \texttt{ADS1299Plus::begin()}. En esta fase se inicializa la comunicación SPI, se aplica un reinicio del dispositivo, se detiene cualquier conversión que pudiera estar activa y se fuerza el modo \texttt{SDATAC}. Este modo permite escribir y leer registros con seguridad antes de iniciar la adquisición continua. Finalmente, se lee el registro de identificación para comprobar que el circuito detectado corresponde a la familia ADS1299 y a la variante de cuatro canales usada en el prototipo.

Después de esta inicialización, \texttt{configureDefaults()} carga la configuración base del convertidor. Esta función programa los registros que fijan la tasa de datos, la referencia interna, el estado de cada canal, la derivación de BIAS y la detección \textit{lead-off}. Por ejemplo, \texttt{CONFIG1} define la frecuencia de muestreo, \texttt{CONFIG3} habilita la referencia interna y el bloque BIAS, \texttt{CHnSET} configura cada canal y su ganacia, y \texttt{BIAS\_SENSP}/\texttt{BIAS\_SENSN} seleccionan qué entradas participan en la señal de BIAS. De esta forma, el programa principal no contiene valores binarios dispersos, sino una configuración agrupada y más fácil de revisar.


En la versión final del prototipo se utiliza la configruación, resumidoa en la Tabla~\ref{tab:impl_ads1299_modo_final}.
\begin{table}[H]
    \centering
    \caption{Configuración final del ADS1299 en el modo de adquisición empleado.}
    \label{tab:impl_ads1299_modo_final}
    \begin{tabular}{p{0.22\textwidth}p{0.33\textwidth}p{0.34\textwidth}}
        \hline
        \textbf{Elemento} & \textbf{Configuración final} & \textbf{Papel en la adquisición} \\
        \hline
        CH1 & Activo, ganancia 24, entrada normal & Canal EEG principal usado por el procesamiento y la sonificación. \\
        CH2--CH4 & Apagados/cortocircuitados en el modo final & Conservados en el payload por contrato, sin interpretarse como EEG activo principal. \\
        BIAS & Derivado de CH1P + CH1N & Referencia de modo común limitada al canal usado. \\
        \textit{Lead-off} & Desactivado & Evita inyección diagnóstica durante la captura final. \\
        Test interno & Desactivado & No se mezcla señal de prueba con la señal real. \\
        Tasa de muestreo & \SI{250}{\hertz} & Frecuencia común entre firmware y backend Python. \\
        \hline
    \end{tabular}
\end{table}

Esta distinción evita una ambigüedad importante: el sistema mantiene un formato de datos de cuatro canales porque el ADS1299-4, el firmware, las capturas y el backend comparten ese contrato. Sin embargo, la validación experimental final se centra en CH1 como canal fisiológico principal.

\subsubsection{Sincronización temporal mediante DRDY y RDATAC}

Una vez configurado el ADS1299, el firmware inicia la adquisición manteniendo activa la señal \texttt{START} y enviando el comando \texttt{RDATAC}. En este modo, el convertidor genera frames de forma continua. La señal \texttt{DRDY}, activa a nivel bajo, indica que una nueva conversión está lista para ser leída.

El firmware no genera muestras por temporización software. Sigue el ritmo impuesto por el ADS1299. Cuando \texttt{DRDY} baja, una interrupción muy breve incrementa el contador \texttt{drdy\_count}. La lectura SPI no se realiza dentro de la interrupción, porque una ISR debe ser corta y no conviene ejecutar operaciones complicadas. Después, el bucle principal consulta el contador y, si existe una conversión pendiente, llama a \texttt{ADS1299Plus::readFrameRDATAC()}.

Cada frame RDATAC leído en esta configuración ocupa \SI{15}{\byte}: tres bytes de estado y cuatro canales de tres bytes cada uno. El driver baja \texttt{CS}, transfiere quince bytes por SPI enviando \texttt{0x00} como byte de relleno y vuelve a subir \texttt{CS}. El resultado se separa en un campo \texttt{status24} y cuatro muestras de canal.

La Figura~\ref{fig:impl_secuencia_drdy_rdatac} resume esta secuencia. La clave es que \texttt{DRDY} solo marca disponibilidad de dato; la transformación de la muestra se realiza después en el bucle principal.

\begin{figure}[H]
    \centering
    \IfFileExists{images/uml/impl_secuencia_drdy_rdatac.pdf}{%
        \includegraphics[width=0.90\textwidth]{images/uml/impl_secuencia_drdy_rdatac.pdf}
    }{%
        \fbox{\parbox{0.9\textwidth}{\centering Diagrama pendiente de exportar desde Mermaid: \texttt{images/uml/impl\_secuencia\_drdy\_rdatac.mmd}}}%
    }
    \caption{Secuencia simplificada de adquisición. La interrupción de \texttt{DRDY} solo marca la disponibilidad de una muestra; la lectura SPI, el filtrado y la publicación del bloque se realizan posteriormente en el bucle principal.}
    \label{fig:impl_secuencia_drdy_rdatac}
\end{figure}

\subsubsection{Reconstrucción, escalado y filtrado de muestras}

Cada frame RDATAC comienza con un campo de estado de 24 bits. El firmware valida que los bits más significativos coincidan con el prefijo esperado \texttt{0xC00000}. Esta comprobación permite detectar frames incoherentes o problemas de sincronía en la comunicación SPI.

Después del estado aparecen las muestras de canal. Cada canal se transmite como tres bytes en orden \textit{MSB first}. El driver reconstruye esos tres bytes mediante \texttt{unpack24()} y realiza la extensión de signo a 32 bits. Este paso es necesario porque el ADS1299 entrega valores con signo; sin esta extensión, una muestra negativa podría interpretarse como un valor positivo de gran magnitud.

Una vez reconstruido el valor digital de cada canal, el firmware aplica el factor de conversión \texttt{LSB\_V}. En la configuración utilizada, este factor vale \SI{2.235e-8}{\volt} por unidad digital del ADC. Por tanto, la conversión inicial a tensión se expresa como:

\[
    v = c_{\mathrm{ADC}} \cdot \texttt{LSB\_V}
\]

donde \(c_{\mathrm{ADC}}\) es el valor digital con signo reconstruido a partir del frame RDATAC.

La señal resultante se mantiene en voltios para aplicar el acondicionamiento digital: eliminación de corriente continua con filtro paso alto de \SI{0.5}{\hertz}, filtro notch de \SI{50}{\hertz} y paso bajo de \SI{40}{\hertz}. Tras el filtrado, la muestra se convierte a microvoltios enteros mediante \texttt{volts\_to\_uV\_i32()}, de forma que el backend recibe valores en microvoltios y no valores digitales crudos del ADC.

La Figura~\ref{fig:impl_frame_rdatac_a_bloque_uv} muestra la transformación entre el frame nativo del ADS1299 y el bloque que finalmente recibe Python. La figura ayuda a distinguir dos niveles: un frame RDATAC representa una muestra multicanal; un bloque \texttt{eeg\_block\_uV} agrupa ocho frames ya validados, escalados y filtrados.

\begin{figure}[H]
    \centering
    \IfFileExists{images/uml/impl_frame_rdatac_a_bloque_uv.pdf}{%
        \includegraphics[width=0.80\textwidth]{images/uml/impl_frame_rdatac_a_bloque_uv.pdf}
    }{%
        \fbox{\parbox{0.9\textwidth}{\centering Diagrama pendiente de exportar desde Mermaid: \texttt{images/uml/impl\_frame\_rdatac\_a\_bloque\_uv.mmd}}}%
    }
    \caption{Transformación entre un frame RDATAC del ADS1299 y el bloque \texttt{eeg\_block\_uV}. Cada frame individual se valida, se reconstruye, se escala y se filtra; después se acumulan ocho muestras para formar el bloque enviado a Python.}
    \label{fig:impl_frame_rdatac_a_bloque_uv}
\end{figure}

\subsubsection{Agrupación y envío de bloques hacia Linux}

El firmware no publica una notificación por cada muestra individual. En su lugar, acumula ocho muestras filtradas en la estructura \texttt{EegBlockUV}, definida en \texttt{streaming.h}. Con una frecuencia de muestreo de \SI{250}{\hertz}, cada bloque representa \SI{32}{\milli\second} de señal.

Cada bloque contiene una cabecera con \texttt{block\_idx}, \texttt{first\_sample\_idx} y \texttt{sample\_count}. Después incluye, para cada una de las ocho muestras, el \texttt{status} del ADS1299 y los cuatro canales expresados en microvoltios. Esta estructura define el contrato entre el firmware y Python.

\begin{table}[H]
    \centering
    \caption{Relación temporal entre muestra, bloque y transporte en la adquisición.}
    \label{tab:impl_tiempos_adquisicion}
    \begin{tabular}{p{0.28\textwidth}p{0.25\textwidth}p{0.36\textwidth}}
        \hline
        \textbf{Elemento} & \textbf{Valor} & \textbf{Significado} \\
        \hline
        Frecuencia de muestreo & \SI{250}{\hertz} & Una muestra nueva aproximadamente cada \SI{4}{\milli\second}. \\
        Tamaño de bloque & 8 muestras & Agrupación mínima antes de publicar por Bridge. \\
        Periodo nominal de bloque & \SI{32}{\milli\second} & Intervalo temporal representado por cada \texttt{eeg\_block\_uV}. \\
        Canales por muestra & 4 & Contrato estructural común a firmware, backend y capturas. \\
        \hline
    \end{tabular}
\end{table}

Cuando un bloque se completa, se introduce en una cola circular \texttt{TxBlockRing}. Esta cola desacopla la adquisición de la publicación por Bridge y permite que el firmware siga adquiriendo aunque el envío de un bloque tarde algo más que una iteración del bucle. La publicación se realiza mediante \texttt{Bridge.notify("eeg\_block\_uV", ...)}, enviando la cabecera y las ocho muestras del bloque.

La Figura~\ref{fig:impl_payload_eeg_block_uv} representa la estructura lógica del payload que se transmite hacia Linux.

\begin{figure}[H]
    \centering
    \IfFileExists{images/uml/impl_payload_eeg_block_uv.pdf}{%
        \includegraphics[width=0.92\textwidth]{images/uml/impl_payload_eeg_block_uv.pdf}
    }{%
        \fbox{\parbox{0.9\textwidth}{\centering Diagrama pendiente de exportar desde Mermaid: \texttt{images/uml/impl\_payload\_eeg\_block\_uv.mmd}}}%
    }
    \caption{Estructura lógica del payload \texttt{eeg\_block\_uV}. El bloque contiene cabecera temporal y ocho muestras; cada muestra conserva el estado ADS1299 y los cuatro canales expresados en microvoltios.}
    \label{fig:impl_payload_eeg_block_uv}
\end{figure}

\subsubsection{Handlers de control y diagnóstico}

Además de la adquisición continua, el firmware expone varios handlers a través de Bridge. Estos handlers permiten consultar el estado del sistema, modificar parámetros de visualización o generación y enviar bytes MIDI hacia la UART física. La Tabla~\ref{tab:impl_handlers_firmware} resume los handlers principales.

\begin{table}[H]
    \centering
    \caption{Handlers Bridge implementados en firmware.}
    \label{tab:impl_handlers_firmware}
    \begin{tabular}{p{0.25\textwidth}p{0.62\textwidth}}
        \hline
        \textbf{Handler} & \textbf{Función} \\
        \hline
        \texttt{status} & Devuelve contadores de adquisición, errores SPI, estado de DRDY y métricas básicas de ejecución. \\
        \texttt{fs} & Permite consultar o ajustar la frecuencia nominal usada por el firmware. \\
        \texttt{rgb} & Actualiza parámetros de visualización LED asociados al estado del sistema. \\
        \texttt{midi\_bytes} & Recibe desde Python un mensaje MIDI ya codificado en bytes y lo escribe por la UART MIDI física. \\
        \hline
    \end{tabular}
\end{table}

En la ruta MIDI, el firmware actúa como pasarela física. El backend decide qué eventos musicales generar, los convierte en mensajes MIDI y llama al handler \texttt{midi\_bytes}. El firmware no interpreta musicalmente esos bytes: únicamente valida su longitud básica y los transmite por \texttt{Serial1}/D1 a \SI{31250}{baudios}, con la configuración eléctrica necesaria para el conector MIDI OUT.

Para facilitar la validación experimental, el firmware mantiene contadores acumulados. Entre ellos se encuentran el número de muestras procesadas, bloques enviados, errores SPI, errores de estado, desbordamientos de buffer y mensajes MIDI emitidos. Estos valores permiten detectar pérdidas de adquisición, problemas de sincronización o fallos en la salida física durante pruebas prolongadas.

\subsubsection{Papel del firmware dentro de la cadena completa}

La implementación firmware queda, por tanto, limitada a tres responsabilidades principales: configurar y leer el ADS1299, publicar bloques EEG calibrados hacia Python y ofrecer una salida física para los bytes MIDI calculados en el backend. Esta separación es importante para la arquitectura general del proyecto, porque evita mezclar procesamiento espectral o reglas musicales dentro del microcontrolador.

En términos de flujo de datos, el firmware transforma:

\[
    \text{frame ADS1299} \rightarrow \text{muestra escalada y filtrada} \rightarrow \text{bloque } \texttt{eeg\_block\_uV}
\]

y, en sentido inverso para la salida sonora:

\[
    \texttt{midi\_bytes} \rightarrow \text{UART MIDI} \rightarrow \text{MIDI OUT físico}
\]

De este modo, la parte de bajo nivel queda acotada y verificable, mientras que las decisiones de análisis EEG y sonificación permanecen en Python, donde son más fáciles de inspeccionar, registrar y modificar durante la validación del prototipo.

\subsection{Implementación del backend Python}

El backend Python constituye la capa intermedia entre la adquisición EEG realizada por el firmware y las salidas de alto nivel del sistema, como la sonificación, la interfaz web y la matriz LED. Su función principal es recibir bloques EEG ya escalados, mantener una ventana temporal de señal, calcular características espectrales y métricas de calidad, y exponer un estado vivo del sistema para el resto de módulos.

A diferencia del firmware, que debe ejecutarse cerca del hardware y responder al ritmo impuesto por el ADS1299, el backend opera sobre bloques ya agrupados. Esto permite separar la temporización estricta de adquisición del análisis de señal y de la generación musical. La consecuencia práctica es que Python no interpreta directamente frames SPI ni cuentas ADC: recibe una secuencia de bloques \texttt{eeg\_block\_uV} en microvoltios y los procesa mediante módulos especializados.

La Figura~\ref{fig:impl_backend_flujo_general} resume la ruta implementada desde la recepción de un bloque EEG hasta la publicación de estado para el resto del sistema. El diagrama se centra en el backend y sirve como puente entre la implementación del firmware, descrita en la sección anterior, y la sonificación, descrita posteriormente.

\begin{figure}[H]
    \centering
    \IfFileExists{images/uml/impl_backend_flujo_general.pdf}{%
        \includegraphics[width=0.95\textwidth]{images/uml/impl_backend_flujo_general.pdf}
    }{%
        \fbox{\parbox{0.9\textwidth}{\centering Diagrama pendiente de exportar desde Mermaid: \texttt{images/uml/impl\_backend\_flujo\_general.mmd}}}%
    }
    \caption{Flujo general del backend Python. Los bloques EEG recibidos por Bridge se validan, se insertan en el buffer temporal, se convierten en características espectrales y se publican como estado de sistema para sonificación, WebUI y validación.}
    \label{fig:impl_backend_flujo_general}
\end{figure}

\subsubsection{Recepción de bloques EEG}

La recepción de datos se implementa en la clase \texttt{EEGReceiver}. Este módulo actúa como punto de entrada del backend para las notificaciones emitidas por el firmware mediante \texttt{Bridge.notify}. Su responsabilidad es transformar el payload textual del evento \texttt{eeg\_block\_uV} en una estructura interna validada, sin realizar todavía procesamiento espectral ni decisiones musicales.

Cada notificación contiene un bloque de ocho muestras, con índice de bloque, primer índice de muestra, número de muestras y, para cada muestra, estado ADS1299 y valores de cuatro canales en microvoltios. El receptor comprueba que el payload tenga una estructura coherente, que el número de muestras coincida con el declarado y que cada muestra incluya todos los canales esperados. Cuando el bloque es válido, lo inserta en una cola interna para que el backend pueda procesarlo en su ciclo principal.

La Figura~\ref{fig:impl_backend_recepcion_bloque} muestra esta conversión entre evento Bridge y bloque Python validado. La figura distingue entre el formato de transporte, todavía textual y orientado a comunicación, y la representación interna usada por el backend.

\begin{figure}[H]
    \centering
    \IfFileExists{images/uml/impl_backend_recepcion_bloque.pdf}{%
        \includegraphics[width=0.90\textwidth]{images/uml/impl_backend_recepcion_bloque.pdf}
    }{%
        \fbox{\parbox{0.9\textwidth}{\centering Diagrama pendiente de exportar desde Mermaid: \texttt{images/uml/impl\_backend\_recepcion\_bloque.mmd}}}%
    }
    \caption{Recepción de un bloque EEG en el backend. El payload publicado por el firmware se parsea, se valida y se almacena como bloque listo para insertarse en el buffer temporal.}
    \label{fig:impl_backend_recepcion_bloque}
\end{figure}

El receptor también mantiene métricas de transporte. Entre ellas se encuentran los bloques recibidos, bloques válidos, bloques descartados, errores de parseo, saltos de índice y edad del último bloque. Estas métricas permiten detectar fallos de comunicación entre firmware y Python y se incorporan posteriormente al estado observado del sistema.

\subsubsection{Buffer temporal de señal}

Una vez validados, los bloques pasan al procesador de señal \texttt{EEGSignalProcessor}. Este módulo mantiene una ventana temporal circular de la señal reciente y ofrece los métodos necesarios para calcular características. La señal se conserva en unidades físicas, ya escalada a microvoltios por el firmware, por lo que el backend no necesita aplicar factores de conversión del ADC.

El buffer temporal separa el ritmo de llegada de bloques del ritmo de cálculo de características. Los bloques llegan con la cadencia de adquisición, mientras que el cálculo espectral se realiza cuando existe suficiente señal acumulada y se cumple el salto configurado. En la configuración usada, el sistema trabaja con una ventana de características de \SI{4}{\second} y un salto de 64 muestras. A \SI{250}{\hertz}, esto implica una actualización aproximada de características cada \SI{256}{\milli\second}.

La Figura~\ref{fig:impl_backend_buffer_temporal} resume esta diferencia entre bloques de entrada, buffer circular y ventana usada para cálculo de características.

\begin{figure}[H]
    \centering
    \IfFileExists{images/uml/impl_backend_buffer_temporal.pdf}{%
        \includegraphics[width=0.90\textwidth]{images/uml/impl_backend_buffer_temporal.pdf}
    }{%
        \fbox{\parbox{0.9\textwidth}{\centering Diagrama pendiente de exportar desde Mermaid: \texttt{images/uml/impl\_backend\_buffer\_temporal.mmd}}}%
    }
    \caption{Relación entre recepción de bloques y ventana de análisis. Los bloques de ocho muestras alimentan un buffer circular; cuando existe señal suficiente, el backend extrae una ventana de \SI{4}{\second} para calcular características espectrales.}
    \label{fig:impl_backend_buffer_temporal}
\end{figure}

\texttt{EEGSignalProcessor} mantiene además contadores de muestras, marcas temporales y diagnósticos de estado. Estos datos permiten saber si el buffer está listo, cuánta señal útil existe, cuándo se calculó la última ventana y si la entrada se ha quedado congelada o sin datos recientes.

\subsubsection{Extracción espectral de características}

Cuando el buffer dispone de señal suficiente, el backend calcula las características mediante el método \texttt{compute\_live\_features()}. Esta función constituye el punto en el que la señal temporal se resume en variables útiles para la sonificación y para la observación del sistema. En la versión final, el análisis principal se realiza sobre CH1, que es el canal EEG activo utilizado en la validación experimental.

Antes del cálculo espectral, la ventana se centra y se trata para reducir la influencia de desplazamientos de continua. Después se estima la densidad espectral de potencia mediante el método multitaper explicado en los conceptos previos. A partir de esta PSD se integran las potencias relativas de las bandas \(\delta\), \(\theta\), \(\alpha\), \(\beta\) y \(\gamma\), se localiza la frecuencia pico y se calcula el RMS de la ventana.

La función devuelve un diccionario de características. Entre los campos principales se encuentran \texttt{rms}, \texttt{peak\_freq}, \texttt{bandpower}, \texttt{band\_peaks}, \texttt{dominant\_band}, \texttt{alpha\_beta\_ratio} y métricas auxiliares de calidad de señal. La Tabla~\ref{tab:impl_backend_features} resume los grupos de información entregados por el backend al resto del sistema.

\begin{table}[H]
    \centering
    \caption{Grupos de características calculadas por el backend.}
    \label{tab:impl_backend_features}
    \begin{tabular}{p{0.28\textwidth}p{0.60\textwidth}}
        \hline
        \textbf{Grupo} & \textbf{Contenido} \\
        \hline
        Amplitud temporal & RMS, RMS en microvoltios, percentiles, valores mínimo y máximo de la ventana. \\
        Espectro & Frecuencia pico, potencia total, potencias por banda y picos por banda. \\
        Bandas EEG & Potencias relativas de \(\delta\), \(\theta\), \(\alpha\), \(\beta\), \(\gamma\), banda dominante y relación alfa/beta. \\
        Diagnósticos & Estado del buffer, edad del último bloque, congelación de entrada, saturación o baja energía. \\
        Salida para sonificación & Diccionario estable usado por \texttt{SonificationFeatureAdapter} como entrada de la capa musical. \\
        \hline
    \end{tabular}
\end{table}

El cálculo de características se mantiene separado de la sonificación. El backend no decide en esta fase qué nota tocar ni qué mensaje MIDI emitir; únicamente proporciona una representación compacta y actualizada de la ventana EEG. Esta separación permite reutilizar las características para la WebUI, capturas, métricas de calidad y generación musical sin duplicar el análisis espectral.

\subsubsection{Evaluación de calidad y estado vivo del backend}

Además de calcular características, el backend evalúa si la ventana disponible es fiable para modular la salida musical. Esta evaluación se realiza mediante \texttt{compute\_spectral\_quality()}, que recibe las características y diagnósticos asociados a la entrada. Su resultado no sustituye a las características EEG, sino que añade una capa de decisión sobre su uso.

El estado de calidad se expresa mediante un \textit{score}, un estado textual y un factor de puerta o \texttt{gate\_factor}. Cuando la señal es suficientemente estable, las características pueden usarse de forma normal. Si se detectan artefactos, congelación de entrada, energía insuficiente o ausencia de datos recientes, el sistema reduce o bloquea su influencia sobre la sonificación. Esta decisión evita que una ventana no fiable produzca una respuesta musical engañosa.

La Figura~\ref{fig:impl_backend_estado_vivo} muestra cómo se combinan bloques recibidos, buffer, características, calidad y snapshot de estado. El snapshot no es una nueva etapa de cálculo EEG, sino una vista agregada para observación, interfaz y validación.

\begin{figure}[H]
    \centering
    \IfFileExists{images/uml/impl_backend_estado_vivo.pdf}{%
        \includegraphics[width=0.92\textwidth]{images/uml/impl_backend_estado_vivo.pdf}
    }{%
        \fbox{\parbox{0.9\textwidth}{\centering Diagrama pendiente de exportar desde Mermaid: \texttt{images/uml/impl\_backend\_estado\_vivo.mmd}}}%
    }
    \caption{Construcción del estado vivo del backend. Las métricas de recepción, el estado del buffer, las características espectrales y la calidad de la ventana se agregan en un snapshot observable para interfaz, validación y sonificación.}
    \label{fig:impl_backend_estado_vivo}
\end{figure}

El snapshot incluye información de adquisición, estado del buffer, última ventana analizada, características espectrales, calidad, errores y métricas de transporte. Esta estructura permite que otros módulos consulten el estado del sistema sin acoplarse directamente al receptor ni al procesador de señal.

\subsubsection{Relación del backend con el resto del sistema}

La salida del backend se utiliza en tres rutas principales. La primera es la sonificación, que consume características y calidad para generar controles musicales. La segunda es la interfaz web, que muestra el estado del sistema y permite modificar ciertos parámetros. La tercera es la validación experimental, donde los snapshots y capturas permiten reconstruir qué estaba ocurriendo durante una prueba.

En términos de responsabilidad, el backend actúa como frontera entre adquisición y decisión musical:

\[
    \texttt{eeg\_block\_uV} \rightarrow \text{buffer temporal} \rightarrow \text{features + calidad} \rightarrow \text{estado vivo}
\]

Esta ruta deja preparada la información necesaria para la siguiente capa. A partir de las características calculadas, la sonificación ya no necesita conocer detalles del ADS1299, del formato del bloque o del cálculo multitaper: únicamente recibe un conjunto estable de variables que puede convertir en parámetros musicales.

\subsection{Implementación de la sonificación}

La sonificación constituye la tercera capa ejecutable de la cadena EEG--MIDI. En las secciones anteriores se ha descrito cómo el firmware adquiere la señal EEG real y cómo el backend Python la transforma en una ventana reciente de características espectrales y métricas de calidad. A partir de ese punto, la sonificación no vuelve a trabajar con muestras crudas ni recalcula la PSD: utiliza el estado EEG ya resumido por el backend para generar música en tiempo real.

La función de esta capa es transformar una representación compacta de la señal en controles musicales. Estos controles son valores normalizados que indican, por ejemplo, cuánta actividad, tensión, densidad rítmica, estabilidad, registro o dinámica debe tener la salida musical. Por tanto, el sistema no realiza una traducción literal del tipo ``el valor de la potencia de una banda EEG equivale a una nota concreta''. En su lugar, las bandas EEG, el RMS, la frecuencia pico y la calidad de la ventana modifican parámetros musicales de alto nivel.

La generación musical mantiene separadas dos fuentes de información. Por un lado, el EEG aporta controles variables en el tiempo. Por otro lado, la configuración musical fija el marco sobre el que se genera la música: escala, nota raíz, nota principal, canal MIDI, programa instrumental y duración del compás. Esta separación evita interpretar la sonificación como una lectura clínica directa de la señal y permite entenderla como un mapeo algorítmico entre características EEG y comportamiento musical.

El flujo completo puede resumirse así: las características EEG y el estado de calidad se adaptan mediante \texttt{SonificationFeatureAdapter}; el resultado se almacena en una estructura \texttt{SonificationFeatures}; después se construye un segmento musical, se genera un compás, se crean notas y finalmente esas notas se programan como eventos MIDI. En la última etapa, los eventos se convierten en bytes MIDI y se envían al firmware mediante \texttt{Bridge.call("midi\_bytes")}, desde donde salen por la UART asociada al conector MIDI físico.

La Figura~\ref{fig:impl_sonificacion_flujo} resume esta ruta. Debe leerse como continuación directa de la sección 6.2: el bloque de procesamiento ya ha calculado las características y la calidad de la ventana, y la sonificación empieza cuando ese estado interno se convierte en controles musicales.

\begin{figure}[H]
    \centering
    \IfFileExists{images/uml/impl_sonificacion_flujo.png}{%
        \includegraphics[width=0.98\textwidth,height=0.42\textheight,keepaspectratio]{images/uml/impl_sonificacion_flujo.png}
    }{%
        \IfFileExists{images/uml/impl_sonificacion_flujo.pdf}{%
            \includegraphics[width=0.98\textwidth,height=0.42\textheight,keepaspectratio]{images/uml/impl_sonificacion_flujo.pdf}
        }{%
            \fbox{\parbox{0.9\textwidth}{\centering Diagrama pendiente de exportar desde Mermaid: \texttt{images/uml/impl\_sonificacion\_flujo.mmd}}}%
        }%
    }
    \caption{Flujo principal de sonificación implementado. El estado EEG calculado por el backend se transforma primero en controles musicales; después se genera una estructura musical formada por segmento, compás y notas; finalmente, las notas se programan como eventos MIDI, se convierten a bytes y se envían al firmware mediante \texttt{midi\_bytes}.}
    \label{fig:impl_sonificacion_flujo}
\end{figure}

\subsubsection{Segmentación EEG}

En esta sección, el término segmentación no se refiere a una segmentación \textit{offline} de registros EEG ni a la detección clínica de estados cerebrales. La implementación trabaja sobre la ventana viva ya calculada por el backend: CH1 se analiza en ventanas de \SI{4}{\second}, actualizadas cada 64 muestras, y el resultado de esa ventana se entrega a la sonificación como un diccionario de características. Por tanto, el primer paso de 6.3 consiste en transformar una ventana EEG ya resumida en controles musicales, no en volver a procesar la señal temporal.

El módulo \texttt{sonification\_features.py} concentra esta frontera. Su documentación interna establece explícitamente que no calcula DSP, no importa \texttt{EEGSignalProcessor} y no depende de \texttt{BackendService}; recibe el diccionario ya calculado por \texttt{compute\_live\_features()} y lo convierte en una estructura \texttt{SonificationFeatures}. Esta decisión mantiene la separación de responsabilidades: el cálculo espectral pertenece a 6.2, mientras que 6.3 empieza con una representación musicalmente utilizable de esas características.

La clase \texttt{SonificationFeatures} agrupa controles normalizados en el intervalo \([0,1]\). Estos controles son reportables en el \textit{snapshot}, en la WebUI, en capturas y en la memoria del TFG. Los principales son \texttt{alpha\_drive}, \texttt{beta\_gamma\_drive}, \texttt{rms\_beta\_activity}, \texttt{band\_driven\_density}, \texttt{spectral\_register}, \texttt{alpha\_stability}, \texttt{rms\_band\_velocity} y \texttt{band\_note\_probability}. El objeto conserva además información trazable, como RMS en voltios y microvoltios, frecuencia pico, banda dominante y potencias por banda.

La Tabla~\ref{tab:impl_sonificacion_controles} resume el significado de estos controles dentro del motor musical. Esta tabla no debe interpretarse como una tabla clínica de bandas EEG, sino como una descripción del mapeo algorítmico utilizado para generar música.

\begin{table}[H]
    \centering
    \caption{Controles principales usados por la sonificación.}
    \label{tab:impl_sonificacion_controles}
    \renewcommand{\arraystretch}{1.15}
    \begin{tabular}{p{0.23\textwidth}p{0.25\textwidth}p{0.28\textwidth}p{0.14\textwidth}}
        \hline
        \textbf{Control} & \textbf{Origen EEG} & \textbf{Interpretación musical} & \textbf{Rango/efecto} \\
        \hline
        \texttt{alpha\_drive} & Relación relativa entre alfa y beta & Tendencia a estabilidad o reposo musical & 0--1 \\
        \texttt{beta\_gamma\_drive} & Peso de beta y gamma & Tensión, síncopa o actividad rápida & 0--1 \\
        \texttt{rms\_beta\_activity} & RMS normalizado combinado con beta/gamma & Actividad global del compás & 0--1 \\
        \texttt{band\_driven\_density} & Densidad derivada de bandas rápidas y RMS & Cantidad de eventos rítmicos & 0--1 \\
        \texttt{spectral\_register} & Frecuencia pico y peso de bandas rápidas & Desplazamiento hacia registro grave/agudo & 0--1 \\
        \texttt{alpha\_stability} & Alfa, theta y nivel RMS & Estabilidad armónica relativa & 0--1 \\
        \texttt{rms\_band\_velocity} & RMS y bandas rápidas & Dinámica o \textit{velocity} MIDI & 0--1 \\
        \texttt{band\_note\_probability} & Densidad espectral musicalizada & Probabilidad de activar notas & 0--1 \\
        \texttt{quality\_gate} & Estado de calidad espectral & Atenuación o bloqueo musical & 0--1 \\
        \hline
    \end{tabular}
\end{table}

El \texttt{quality\_gate} se incorpora a la sonificación antes de generar nuevo material musical. Si la calidad de la ventana es alta, los controles se aplican con normalidad; si la ventana presenta artefactos o baja fiabilidad, los controles sensibles se atenúan. En particular, la implementación reduce la actividad, la densidad, la velocidad y la probabilidad de nota cuando el factor de puerta disminuye. Si no existen características válidas o el \textit{score} cae por debajo del umbral de uso, la sonificación deja de generar nuevo material útil hasta recibir una ventana posterior válida.

La Figura~\ref{fig:impl_sonificacion_quality_gate} muestra esta relación entre características, diagnósticos, calidad y controles musicales. La finalidad de la figura es aclarar que la calidad no modifica la adquisición ni recalcula el DSP, sino que actúa como una puerta entre las características EEG y el motor musical.

\begin{figure}[H]
    \centering
    \IfFileExists{images/uml/impl_sonificacion_quality_gate.png}{%
        \includegraphics[width=0.95\textwidth,height=0.42\textheight,keepaspectratio]{images/uml/impl_sonificacion_quality_gate.png}
    }{%
        \IfFileExists{images/uml/impl_sonificacion_quality_gate.pdf}{%
            \includegraphics[width=0.95\textwidth,height=0.42\textheight,keepaspectratio]{images/uml/impl_sonificacion_quality_gate.pdf}
        }{%
            \fbox{\parbox{0.9\textwidth}{\centering Diagrama pendiente de exportar desde Mermaid: \texttt{images/uml/impl\_sonificacion\_quality\_gate.mmd}}}%
        }%
    }
    \caption{Papel del \textit{quality gate} en la sonificación. La calidad de la ventana decide si los controles se aplican plenamente, se atenúan o se evita generar nuevo material musical a partir de una ventana no fiable.}
    \label{fig:impl_sonificacion_quality_gate}
\end{figure}

Además de aplicar la puerta de calidad, \texttt{SonificationFeatureAdapter} introduce memoria temporal. El RMS se normaliza respecto a una línea base adaptativa lenta y varios controles se suavizan mediante una media exponencial. Esta decisión reduce cambios bruscos entre ventanas consecutivas y evita que pequeñas fluctuaciones espectrales se traduzcan en saltos musicales demasiado abruptos. El suavizado no oculta que una ventana sea mala: los metadatos de calidad se conservan y el estado \texttt{valid} sigue condicionando la generación posterior.

\subsubsection{Construcción de segmentos musicales}

Una vez obtenidos los controles de sonificación, el backend construye un estado musical intermedio mediante \texttt{MusicSegmentBuilder}. Esta etapa recibe los controles ya suavizados, la escala activa, la nota principal y el tiempo musical actual. El resultado es un objeto \texttt{MusicSegment}, que no contiene todavía notas individuales, sino la información necesaria para generar el siguiente compás.

La función de \texttt{MusicSegmentBuilder} es traducir los controles continuos de 6.3.1 a campos musicales más directamente utilizables. Por ejemplo, la densidad rítmica continua se convierte en una cadencia discreta, el registro espectral se guarda como pista de registro, la estabilidad pasa a ser un criterio armónico y la dinámica se conserva como factor de \textit{velocity}. De este modo, el sistema separa dos pasos: primero calcula controles a partir del EEG y después organiza esos controles como estado musical.

La Figura~\ref{fig:impl_sonificacion_music_segment_builder} muestra esta reorganización. El diagrama debe leerse de izquierda a derecha: los controles de \texttt{SonificationFeatures} entran en \texttt{MusicSegmentBuilder}; el builder los transforma o copia a campos musicales concretos; y el resultado queda agrupado en \texttt{MusicSegment}, que será la entrada del generador de compases.

\begin{figure}[H]
    \centering
    \IfFileExists{images/uml/impl_sonificacion_music_segment_builder.png}{%
        \includegraphics[width=0.98\textwidth,height=0.48\textheight,keepaspectratio]{images/uml/impl_sonificacion_music_segment_builder.png}
    }{%
        \IfFileExists{images/uml/impl_sonificacion_music_segment_builder.pdf}{%
            \includegraphics[width=0.98\textwidth,height=0.48\textheight,keepaspectratio]{images/uml/impl_sonificacion_music_segment_builder.pdf}
        }{%
            \fbox{\parbox{0.9\textwidth}{\centering Diagrama pendiente de exportar desde Mermaid: \texttt{images/uml/impl\_sonificacion\_music\_segment\_builder.mmd}}}%
        }%
    }
    \caption{Construcción del \texttt{MusicSegment}. Los controles de sonificación se reorganizan en campos musicales: cadencia rítmica, registro, tensión, estabilidad, actividad, dinámica y probabilidad de nota. La escala, la nota principal y la duración proceden de la configuración musical.}
    \label{fig:impl_sonificacion_music_segment_builder}
\end{figure}

La transformación más importante de esta etapa es la conversión de \texttt{band\_driven\_density} en \texttt{rhythm\_cadence}. El valor de densidad es continuo, pero el generador de compases trabaja con tres niveles discretos: \texttt{LOW}, \texttt{MEDIUM} y \texttt{HIGH}. Una cadencia baja tiende a compases con menos notas; una cadencia media produce un patrón intermedio; y una cadencia alta permite compases con más posiciones activas. Para evitar saltos constantes entre estados, el cambio de cadencia se realiza con histéresis: una pequeña variación entre ventanas consecutivas no provoca necesariamente un cambio de ritmo.

La Tabla~\ref{tab:impl_music_segment_mapeo} resume los campos principales de \texttt{MusicSegment}. En la columna central se indica qué entrada usa el builder, evitando expresiones ambiguas como ``derivado de''. La última columna indica para qué se usará ese campo en las etapas posteriores.

\begin{table}[H]
    \centering
    \caption{Mapa de entradas y campos principales de \texttt{MusicSegment}.}
    \label{tab:impl_music_segment_mapeo}
    \footnotesize
    \setlength{\tabcolsep}{3pt}
    \renewcommand{\arraystretch}{1.10}
    \begin{tabular}{p{0.27\textwidth}p{0.31\textwidth}p{0.33\textwidth}}
        \hline
        \textbf{Campo en \texttt{MusicSegment}} & \textbf{Entrada usada por el builder} & \textbf{Uso posterior} \\
        \hline
        \texttt{rhythm\_cadence} &
        \texttt{band\_driven\_density} convertida a \texttt{LOW}, \texttt{MEDIUM} o \texttt{HIGH}. &
        Define la cantidad base de posiciones activas del compás. \\
        \texttt{register\_hint} &
        Valor de \texttt{spectral\_register}. &
        Desplaza el centro melódico hacia registros graves o agudos. \\
        \texttt{tension} &
        Valor de \texttt{beta\_gamma\_drive}. &
        Favorece síncopas, movimiento rítmico y mayor tensión melódica o armónica. \\
        \texttt{harmonic\_stability} &
        Valor de \texttt{alpha\_stability}. &
        Ayuda a elegir acordes más estables o más tensos. \\
        \texttt{activity} &
        Valor de \texttt{rms\_beta\_activity}. &
        Ajusta la presencia general del compás y puede aumentar ligeramente los eventos. \\
        \texttt{velocity\_factor} &
        Valor de \texttt{rms\_band\_velocity}. &
        Modula la intensidad de las notas MIDI. \\
        \texttt{note\_probability} &
        Valor de \texttt{band\_note\_probability}. &
        Ajusta la probabilidad de que una posición rítmica genere una nota. \\
        \texttt{scale} y \texttt{main\_note\_midi} &
        Escala y nota principal configuradas en el sistema. &
        Mantienen el marco tonal y el centro musical de la melodía. \\
        \hline
    \end{tabular}
\end{table}

El segmento musical también almacena el intervalo temporal del siguiente compás. En la ruta en tiempo real, esta duración viene dada por \texttt{MUSIC\_BAR\_SEC}, cuyo valor por defecto es \SI{2}{\second}. Por tanto, cada \texttt{MusicSegment} representa el estado musical que se usará para generar un compás de esa duración, asociado al estado EEG reciente.

Finalmente, el campo \texttt{valid} permite mantener la relación con la calidad de la sonificación. Si el estado anterior no es válido, el backend no utiliza ese resultado para generar nuevo material musical. Así, \texttt{MusicSegment} funciona como una frontera clara entre los controles continuos y la generación de compases: cuando el estado es válido, el objeto se entrega a \texttt{BarGenerator}; cuando no lo es, la generación se bloquea hasta recibir una actualización fiable.

En resumen, \texttt{MusicSegmentBuilder} convierte los controles de sonificación en una descripción musical compacta del siguiente compás. A partir de ese momento, el sistema ya dispone de escala, nota principal, cadencia rítmica, registro, tensión, estabilidad, actividad, dinámica y probabilidad de nota. El siguiente apartado explica cómo \texttt{BarGenerator} usa esta información para construir el compás propiamente dicho.

\subsubsection{Generación de compases}

El siguiente nivel de la jerarquía es el compás. \texttt{BarGenerator} recibe un \texttt{MusicSegment} y produce un objeto \texttt{Bar}. Este objeto contiene el índice del compás, tiempos de inicio y fin, la raíz del acorde, las notas del acorde, un patrón de posiciones activas y una envolvente musical sintética. El compás funciona como puente entre controles continuos y notas individuales.

La implementación usa una cuadrícula de 16 posiciones por compás: cuatro pulsos con cuatro subdivisiones por pulso. Cada posición puede estar activa o en silencio. El primer slot se fuerza como ancla musical para que el compás tenga un punto de referencia estable. El número objetivo de notas depende de la cadencia rítmica \texttt{LOW}, \texttt{MEDIUM} o \texttt{HIGH}, y se ajusta ligeramente con la actividad estimada.

Los pesos de selección de slots combinan reglas musicales y controles EEG. La base favorece el primer pulso, los tiempos principales y ciertas subdivisiones. Después, la actividad incrementa la probabilidad global de eventos, la tensión favorece síncopas u \textit{offbeats}, y \texttt{band\_note\_probability} ajusta la probabilidad general de activar posiciones. Por tanto, el EEG no decide manualmente cada slot, sino que modifica una distribución musical ponderada.

En paralelo, el compás selecciona un acorde diatónico a partir de la escala activa. La estabilidad armónica y la tensión determinan si se tiende a un grado estable, intermedio o de mayor tensión. La implementación utiliza histéresis armónica para evitar saltos continuos entre grados por cambios pequeños de la señal. Esto mantiene una salida más coherente en tiempo real.

La Figura~\ref{fig:impl_sonificacion_segmento_compas_nota} resume la relación entre segmento, compás y notas. La misma figura se utiliza también como guía para la subsección siguiente, porque muestra que las notas se generan después de disponer de escala, acorde, patrón rítmico y controles musicales.

\begin{figure}[H]
    \centering
    \IfFileExists{images/uml/impl_sonificacion_segmento_compas_nota.png}{%
        \includegraphics[width=0.98\textwidth,height=0.46\textheight,keepaspectratio]{images/uml/impl_sonificacion_segmento_compas_nota.png}
    }{%
        \IfFileExists{images/uml/impl_sonificacion_segmento_compas_nota.pdf}{%
            \includegraphics[width=0.98\textwidth,height=0.46\textheight,keepaspectratio]{images/uml/impl_sonificacion_segmento_compas_nota.pdf}
        }{%
            \fbox{\parbox{0.9\textwidth}{\centering Diagrama pendiente de exportar desde Mermaid: \texttt{images/uml/impl\_sonificacion\_segmento\_compas\_nota.mmd}}}%
        }%
    }
    \caption{Jerarquía musical empleada por la sonificación. Los controles EEG suavizados y la configuración musical construyen un segmento; el segmento genera un compás con acorde y patrón rítmico; el generador de notas produce eventos \texttt{NoteEvent}.}
    \label{fig:impl_sonificacion_segmento_compas_nota}
\end{figure}

\subsubsection{Generación de notas}

\texttt{NoteGenerator} convierte el \texttt{MusicSegment} y el \texttt{Bar} en una lista de \texttt{NoteEvent}. Esta estructura representa una nota musical de alto nivel, pero todavía no es un mensaje MIDI. Cada \texttt{NoteEvent} contiene tiempo de inicio, tiempo de fin, altura MIDI, \textit{velocity}, canal y programa. La conversión a \texttt{note\_on} y \texttt{note\_off} se realiza después, en el scheduler MIDI.

La selección de alturas se realiza dentro de la escala configurada. Primero se calcula un centro de registro alrededor de la nota principal. El control \texttt{spectral\_register} desplaza suavemente ese centro hacia registros más graves o más agudos, pero la nota principal sigue actuando como referencia musical. A partir de ese centro se construye un conjunto de notas de la escala dentro de un radio determinado, y se separan notas de acorde y notas de paso.

Las posiciones fuertes del compás tienden a utilizar notas del acorde, mientras que las posiciones intermedias pueden usar notas de paso cuando existen candidatos adecuados. La tensión EEG desplaza el objetivo melódico y puede permitir intervalos algo mayores, mientras que la actividad y la probabilidad influyen en la variedad. Aun así, el generador limita los saltos respecto a la nota anterior para evitar melodías erráticas y mantiene las alturas dentro del rango MIDI permitido.

La dinámica se calcula mediante la función de \textit{velocity}. Esta combina un acento musical base según la posición rítmica, el factor EEG \texttt{rms\_band\_velocity} y la envolvente sintética del compás. El resultado se acota entre una velocidad mínima y máxima para evitar notas demasiado débiles o excesivamente agresivas. En este punto, el EEG afecta a la intensidad y al carácter de la nota, pero no se interpreta como una medida clínica de activación cerebral.

La duración de cada nota se obtiene a partir de la cuadrícula del compás. Una nota dura hasta el siguiente slot activo o hasta el final del compás. Si por seguridad se detectara una duración nula o negativa, se impone una duración mínima. Esta regla es importante para el flujo MIDI posterior, porque cada nota debe tener un inicio y un final definidos para evitar notas colgadas.

En el primer slot, el sistema puede generar el acorde completo como marcador armónico. El backend decide si se reproduce ese acorde completo según un periodo mínimo y según si los controles musicales han cambiado de forma suficiente. Así se evita que cada compás repita un acorde completo de forma innecesaria, manteniendo una textura más ligera cuando el estado EEG no cambia significativamente.

\subsubsection{Escritura y transmisión MIDI}

La escritura MIDI se divide en dos niveles. El primer nivel es musical y trabaja con \texttt{NoteEvent}; el segundo nivel es MIDI y trabaja con eventos temporizados y bytes. Esta separación permite que \texttt{NoteGenerator} no tenga que conocer el transporte físico, y que \texttt{MidiByteTransport} no tenga que conocer las reglas de generación musical.

\texttt{MidiScheduler}, implementado en \texttt{midi\_live.py}, transforma cada \texttt{NoteEvent} en dos eventos MIDI vivos: un \texttt{note\_on} en el inicio de la nota y un \texttt{note\_off} en su final. Además, el backend programa un \texttt{program\_change} una sola vez al inicio para fijar el programa MIDI de salida. Los eventos se almacenan en una cola temporal ordenada por \texttt{due\_time}, usando el reloj monotónico del sistema.

En cada iteración del backend, \texttt{\_pump\_midi()} llama a \texttt{pop\_due\_events(now, lookahead\_sec=MIDI\_LOOKAHEAD\_SEC)}. El \textit{lookahead} de \SI{0.02}{\second} permite extraer eventos ligeramente por adelantado, de modo que el transporte y el firmware dispongan de margen para emitirlos sin bloquear el ciclo principal. El scheduler no envía nada por sí mismo: solo decide qué eventos están vencidos o próximos a vencer.

La conversión a bytes se realiza mediante \texttt{event\_to\_midi\_bytes()}. Un \texttt{note\_on} se codifica con estado \texttt{0x90 | channel}, nota y \textit{velocity}; un \texttt{note\_off} usa \texttt{0x80 | channel}; un \texttt{control\_change} usa \texttt{0xB0 | channel}; y un \texttt{program\_change} usa \texttt{0xC0 | channel}. Todos los canales se limitan a 0--15 y los datos MIDI a 0--127.

\texttt{MidiByteTransport} recibe los eventos vencidos, los convierte en bytes y llama al \textit{handler} del firmware mediante \texttt{Bridge.call("midi\_bytes", n, b0, b1, b2)}. El campo \texttt{n} indica si el mensaje tiene dos o tres bytes válidos. El firmware no decide qué nota tocar ni interpreta musicalmente el mensaje: recibe los bytes y los escribe por la UART asociada a la salida MIDI física.

La salida física validada usa \texttt{Serial1}/D1 a la velocidad MIDI estándar de \SI{31250}{baudios}, con transmisión invertida en el microcontrolador. Esta configuración pertenece a la capa de firmware y electrónica ya descrita previamente; en 6.3 basta con indicar que la ruta Python--Bridge--firmware desemboca en esa UART física y, desde ella, en el conector MIDI OUT del prototipo.

La Figura~\ref{fig:impl_midi_eventos_bytes} resume la formación de mensajes MIDI desde las notas generadas. Esta figura aclara la diferencia entre nota musical, evento MIDI temporizado, bytes MIDI y transporte físico.

\begin{figure}[H]
    \centering
    \IfFileExists{images/uml/impl_midi_eventos_bytes.png}{%
        \includegraphics[width=0.92\textwidth,height=0.62\textheight,keepaspectratio]{images/uml/impl_midi_eventos_bytes.png}
    }{%
        \IfFileExists{images/uml/impl_midi_eventos_bytes.pdf}{%
            \includegraphics[width=0.92\textwidth,height=0.62\textheight,keepaspectratio]{images/uml/impl_midi_eventos_bytes.pdf}
        }{%
            \fbox{\parbox{0.9\textwidth}{\centering Diagrama pendiente de exportar desde Mermaid: \texttt{images/uml/impl\_midi\_eventos\_bytes.mmd}}}%
        }%
    }
    \caption{Formación de mensajes MIDI en la ruta live. Las notas de alto nivel se transforman en eventos \texttt{note\_on}, \texttt{note\_off} y \texttt{program\_change}; después se convierten a bytes y se envían al firmware mediante \texttt{midi\_bytes}.}
    \label{fig:impl_midi_eventos_bytes}
\end{figure}

Como medida de seguridad, el sistema conserva la ruta \texttt{panic}. Esta acción limpia el scheduler y genera mensajes \texttt{All Sound Off} y \texttt{All Notes Off} para todos los canales. Es necesaria porque una ventana de mala calidad puede bloquear la generación de nuevos compases, pero no elimina automáticamente eventos que ya hubieran sido programados desde una ventana anterior. Por tanto, el \textit{panic} forma parte de la robustez del sistema MIDI y debe mantenerse aunque se simplifique la lógica musical.

La salida MIDI puede desactivarse mediante la variable de entorno \texttt{EEG\_MIDI\_LIVE\_ENABLED}. En ese caso, los eventos pueden seguir generándose y contabilizándose, pero el transporte los marca como descartados en lugar de enviarlos al firmware. Esta opción permite validar generación y scheduler sin forzar siempre el envío físico por la UART MIDI.

\subsection{Implementación de la interfaz web}

% Pendiente de redacción propia.

\subsection{Implementación de la matriz LED}

% Pendiente de redacción propia.

\end{document}
